\chapter*{Abstract}

\noindent Rolling-element bearings are small components, but their condition often decides whether a rotating machine can keep running safely. A bearing fault can begin as a local surface defect and then grow into severe vibration, heat, loss of alignment, secondary damage, and unplanned shutdown. Remaining useful life prediction is therefore a central problem in prognostics and health management. This thesis studies bearing remaining useful life prediction using the NASA IMS run-to-failure bearing data and compares a physics-informed neural network against three data-driven baselines.

The original proposal for this work expected a physics-informed model to outperform a purely data-driven model. The completed local experiments gave a more measured result. The proposed DeepXDE physics-informed model performed best in one of the three cross-bearing experiments, but the LSTM baseline gave the lowest average RMSE and highest average $R^2$ across the full experiment set. This does not make the physics-informed approach unimportant. It shows that weak physical regularization, when applied to noisy run-to-failure vibration features, must be designed and weighted carefully. A simple monotonicity and fault-frequency prior can help in some transfer settings, but it may also restrict the network when the degradation path of the test bearing differs from the training bearings.

The final pipeline reads the IMS data locally from extracted folders, removes the dependency on Google Drive, extracts time-domain and envelope spectral features, constructs normalized RUL targets, builds a PCA-based health indicator, and evaluates four models: a feed-forward neural baseline, the proposed DeepXDE physics-informed model, an LSTM sequence model, and a 1D CNN sequence model. The experiments use three train-validation-test splits designed to test cross-bearing and cross-dataset transfer. The best average RMSE was achieved by the LSTM baseline at 0.162, followed by the data-only feed-forward baseline at 0.185, the proposed physics-informed model at 0.223, and the CNN baseline at 0.235. The work contributes a reproducible local project structure, a documented feature and evaluation pipeline, and an honest analysis of where the investigated physics-informed formulation helped and where it did not.

\textbf{Keywords:} remaining useful life, rolling-element bearing, prognostics and health management, physics-informed neural network, IMS bearing dataset, vibration analysis, LSTM, CNN.
